\chapter{L13 -- Criterio di Popov--Belevitch--Hautus per la raggiungibilità}

\section{Obiettivo della lezione}

Nella lezione precedente abbiamo visto che, dopo aver separato il sistema nella sua parte raggiungibile e nella sua parte non raggiungibile, gli autovalori del blocco non raggiungibile non possono essere influenzati dall’ingresso.  
Nasce allora una domanda naturale:

\begin{quote}
\emph{Esiste un criterio puramente algebrico che permetta di capire, guardando \(A\) e \(B\), se un sistema è completamente raggiungibile?}
\end{quote}

La risposta è sì, ed è data dal \textbf{criterio PBH} (Popov--Belevitch--Hautus).

L’idea profonda del criterio è la seguente:
\begin{itemize}
    \item non basta sapere quali sono gli autovalori di \(A\);
    \item bisogna capire se l’ingresso \(B\) riesce davvero ad agire su tutti i modi associati a quegli autovalori;
    \item questo controllo si fa studiando il rango della matrice aumentata
    \[
    \begin{bmatrix}
    A-\lambda I & B
    \end{bmatrix}.
    \]
\end{itemize}

In questa lezione vedremo:
\begin{enumerate}
    \item l’enunciato del criterio PBH;
    \item il suo significato geometrico e modale;
    \item il ruolo della forma di Jordan;
    \item esempi in cui il test si applica in modo molto rapido.
\end{enumerate}

\section{Richiamo: forma standard di raggiungibilità}

Ricordiamo che, se il sistema non è completamente raggiungibile, esiste un cambio di coordinate che porta la coppia \((A,B)\) nella forma standard di raggiungibilità:
\begin{equation}
\tilde A=
\begin{bmatrix}
A_R & A_{RN}\\
0 & A_N
\end{bmatrix},
\qquad
\tilde B=
\begin{bmatrix}
B_R\\
0
\end{bmatrix},
\label{eq:L13_standard_reach}
\end{equation}
dove:
\begin{itemize}
    \item \((A_R,B_R)\) descrive il sottosistema raggiungibile;
    \item \(A_N\) descrive il sottosistema non raggiungibile.
\end{itemize}

Questa scrittura è importantissima, perché ci dice che la parte non raggiungibile evolve da sola e non può essere modificata dall’ingresso.

\begin{figure}[H]
\centering
\begin{tikzpicture}[>=Latex, node distance=3.4cm]
    \node[draw, rounded corners, minimum width=3.4cm, minimum height=1.2cm, fill=blue!8] (zr) {$z_R$};
    \node[draw, rounded corners, minimum width=3.4cm, minimum height=1.2cm, fill=gray!10, right of=zr] (zn) {$z_N$};

    \draw[->, thick] (-2.2,0) -- (zr.west) node[midway, above] {$u$};
    \draw[loop above, thick] (zr) to node[above] {$A_R$} (zr);
    \draw[loop above, thick] (zn) to node[above] {$A_N$} (zn);
    \draw[->, thick] (zr.east) -- (zn.west) node[midway, above] {$A_{RN}$};

    \node[below=0.7cm of zr] {\small parte raggiungibile};
    \node[below=0.7cm of zn] {\small parte non raggiungibile};
\end{tikzpicture}
\caption{Forma standard di raggiungibilità: l’ingresso agisce solo sul sottosistema raggiungibile.}
\end{figure}

\section{Enunciato del criterio PBH}

\begin{lemma}[Criterio PBH per la raggiungibilità]
Sia data la coppia \((A,B)\), con
\[
A\in\R^{n\times n}, \qquad B\in\R^{n\times m}.
\]
La coppia \((A,B)\) è \textbf{completamente raggiungibile} se e solo se
\begin{equation}
\operatorname{rank}
\begin{bmatrix}
A-\lambda I & B
\end{bmatrix}
=n
\qquad \forall \lambda\in\C.
\label{eq:L13_PBH}
\end{equation}
\end{lemma}

\begin{osservazione}
Anche se \(A\) e \(B\) sono reali, il test va formulato per \(\lambda\in\C\), perché gli autovalori di \(A\) possono essere complessi.
\end{osservazione}

\begin{osservazione}
Questo criterio è spesso chiamato anche \emph{criterio PBH di controllabilità}.  
Nel contesto di queste lezioni lo leggiamo come criterio di \textbf{raggiungibilità} della coppia \((A,B)\).
\end{osservazione}

\section{Perché il criterio ha senso}

La matrice
\[
\begin{bmatrix}
A-\lambda I & B
\end{bmatrix}
\]
controlla simultaneamente due cose:
\begin{itemize}
    \item la struttura modale associata a \(\lambda\), tramite \(A-\lambda I\);
    \item la capacità dell’ingresso di agire su quella struttura, tramite \(B\).
\end{itemize}

Se questa matrice perde rango, significa che esiste almeno una direzione modale associata a \(\lambda\) che l’ingresso non riesce a eccitare.

In altre parole:
\begin{quote}
\emph{un autovalore è problematico non solo perché esiste, ma perché può essere associato a un modo che l’ingresso non riesce a raggiungere.}
\end{quote}

\section{Solo gli autovalori sono critici}

Il primo fatto utile è che non serve davvero controllare \eqref{eq:L13_PBH} per tutti i complessi in modo cieco.

\begin{proposizione}
Se \(\lambda\notin\sigma(A)\), allora
\[
\operatorname{rank}
\begin{bmatrix}
A-\lambda I & B
\end{bmatrix}=n.
\]
\end{proposizione}

\begin{proof}
Se \(\lambda\notin\sigma(A)\), allora \(A-\lambda I\) è invertibile, quindi ha già rango \(n\).  
Aggiungere le colonne di \(B\) non può diminuire il rango. Dunque la matrice aumentata ha rango \(n\).
\end{proof}

\begin{osservazione}
Quindi, in pratica, il criterio PBH si controlla solo sugli \textbf{autovalori di \(A\)}.
\end{osservazione}

\section{Dimostrazione del criterio PBH}

\subsection{Necessità}

Supponiamo che il sistema sia completamente raggiungibile.  
Vogliamo dimostrare che allora \eqref{eq:L13_PBH} vale per ogni \(\lambda\in\C\).

Procediamo per assurdo. Supponiamo che esista \(\hat\lambda\in\C\) tale che
\[
\operatorname{rank}
\begin{bmatrix}
A-\hat\lambda I & B
\end{bmatrix}
<n.
\]
Allora esiste un vettore riga non nullo \(w^\top\in\C^{1\times n}\) tale che
\begin{equation}
w^\top
\begin{bmatrix}
A-\hat\lambda I & B
\end{bmatrix}
=0.
\label{eq:L13_left_annihilator}
\end{equation}
Questo equivale a scrivere il sistema
\begin{equation}
\begin{cases}
w^\top(A-\hat\lambda I)=0,\\
w^\top B=0.
\end{cases}
\label{eq:L13_left_eig_system}
\end{equation}
La prima relazione dice che
\[
w^\top A=\hat\lambda\, w^\top,
\]
cioè \(w^\top\) è un \textbf{autovettore sinistro} di \(A\) associato a \(\hat\lambda\).

A questo punto, moltiplicando ripetutamente per \(A\), otteniamo:
\[
w^\top AB=\hat\lambda\, w^\top B=0,
\]
\[
w^\top A^2B=\hat\lambda\, w^\top AB=0,
\]
e in generale
\[
w^\top A^kB=0
\qquad \forall k\ge 0.
\]
Dunque \(w^\top\) annulla tutte le colonne della matrice di raggiungibilità
\[
R=\begin{bmatrix}
B & AB & A^2B & \cdots & A^{n-1}B
\end{bmatrix},
\]
cioè
\[
w^\top R=0.
\]
Ma se il sistema è completamente raggiungibile, allora \(\operatorname{rank}(R)=n\), quindi le colonne di \(R\) generano tutto \(\R^n\), e non può esistere un vettore riga non nullo che le annulli tutte.

Assurdo.  
Quindi necessariamente
\[
\operatorname{rank}
\begin{bmatrix}
A-\lambda I & B
\end{bmatrix}=n
\qquad \forall \lambda\in\C.
\]

\subsection{Sufficienza}

Supponiamo ora che
\[
\operatorname{rank}
\begin{bmatrix}
A-\lambda I & B
\end{bmatrix}=n
\qquad \forall \lambda\in\C.
\]
Vogliamo mostrare che il sistema è completamente raggiungibile.

Procediamo ancora per assurdo. Supponiamo che il sistema \emph{non} sia completamente raggiungibile.  
Allora, per la lezione precedente, esiste un cambio di coordinate che porta il sistema nella forma standard di raggiungibilità \eqref{eq:L13_standard_reach}:
\[
\tilde A=
\begin{bmatrix}
A_R & A_{RN}\\
0 & A_N
\end{bmatrix},
\qquad
\tilde B=
\begin{bmatrix}
B_R\\
0
\end{bmatrix},
\]
con \(\dim(A_N)>0\).

Sia \(\hat\lambda\) un autovalore di \(A_N\).  
Allora esiste un vettore riga non nullo \(\hat x^\top\) tale che
\[
\hat x^\top A_N=\hat\lambda\,\hat x^\top.
\]
Consideriamo allora il vettore riga
\[
\begin{bmatrix}
0^\top & \hat x^\top
\end{bmatrix}.
\]
Moltiplicandolo per la matrice PBH del sistema trasformato otteniamo
\[
\begin{bmatrix}
0^\top & \hat x^\top
\end{bmatrix}
\begin{bmatrix}
\tilde A-\hat\lambda I & \tilde B
\end{bmatrix}
=
\begin{bmatrix}
0^\top & \hat x^\top(A_N-\hat\lambda I) & 0
\end{bmatrix}
=0.
\]
Quindi
\[
\operatorname{rank}
\begin{bmatrix}
\tilde A-\hat\lambda I & \tilde B
\end{bmatrix}<n.
\]
Ma il rango è invariante per similitudine, dunque anche per il sistema originario si avrebbe
\[
\operatorname{rank}
\begin{bmatrix}
A-\hat\lambda I & B
\end{bmatrix}<n,
\]
in contraddizione con l’ipotesi.

Quindi il sistema deve essere completamente raggiungibile.
\qed

\section{Lettura modale del criterio}

Il criterio PBH ci dice che un modo associato all’autovalore \(\lambda\) è \emph{raggiungibile} se e solo se le colonne di \(B\) riescono a rompere la dipendenza lineare eventualmente presente in \(A-\lambda I\).

Dal punto di vista intuitivo:
\begin{itemize}
    \item \(A-\lambda I\) perde rango esattamente quando \(\lambda\) è un autovalore;
    \item allora compaiono direzioni modali “speciali”;
    \item se \(B\) non riesce ad agire in quelle direzioni, l’autovalore corrispondente genera un modo non raggiungibile.
\end{itemize}

Questo collega direttamente il criterio PBH al linguaggio degli autovalori raggiungibili e non raggiungibili visto nella lezione precedente.

\section{Legame con la forma di Jordan}

Per capire ancora meglio il criterio, è utile ragionare in forma di Jordan.

Supponiamo di avere una matrice in forma di Jordan
\[
J=\operatorname{diag}\!\bigl(
J_{\lambda}^{(q_1)},
J_{\lambda}^{(q_2)},
\dots,
J_{\lambda}^{(q_k)},
J_{\mu_1}^{(p_1)},
\dots
\bigr),
\]
dove i primi \(k\) blocchi sono associati allo stesso autovalore \(\lambda\).

\begin{figure}[H]
\centering
\begin{tikzpicture}[node distance=1.7cm]
    \node[draw, minimum width=1.8cm, minimum height=1.8cm] (j1) {$J_\lambda^{(q_1)}$};
    \node[draw, minimum width=1.6cm, minimum height=1.6cm, right=of j1] (j2) {$J_\lambda^{(q_2)}$};
    \node[right=0.9cm of j2] (dots) {$\cdots$};
    \node[draw, minimum width=1.8cm, minimum height=1.8cm, right=0.9cm of dots] (jk) {$J_\lambda^{(q_k)}$};

    \draw[rounded corners, thick]
    ($(j1.north west)+(-0.5,0.5)$) rectangle ($(jk.south east)+(0.5,-0.5)$);

    \node[above=0.4cm of dots] {\(\lambda\)};
\end{tikzpicture}
\caption{Più blocchi di Jordan associati allo stesso autovalore \(\lambda\). Il loro numero è \(m_g(\lambda)\).}
\end{figure}

Allora
\[
J-\lambda I=
\operatorname{diag}\!\bigl(
J_\lambda^{(q_1)}-\lambda I,
J_\lambda^{(q_2)}-\lambda I,
\dots,
J_\lambda^{(q_k)}-\lambda I,
J_{\mu_1}^{(p_1)}-\lambda I,
\dots
\bigr).
\]
Ora:
\begin{itemize}
    \item se \(\mu_i\neq \lambda\), il blocco \(J_{\mu_i}^{(p_i)}-\lambda I\) è invertibile;
    \item se il blocco è \(J_\lambda^{(q)}\), allora
    \[
    J_\lambda^{(q)}-\lambda I
    =
    \begin{bmatrix}
    0 & 1 & 0 & \cdots & 0\\
    0 & 0 & 1 & \cdots & 0\\
    \vdots & \vdots & \ddots & \ddots & \vdots\\
    0 & 0 & \cdots & 0 & 1\\
    0 & 0 & \cdots & 0 & 0
    \end{bmatrix},
    \]
    che ha rango \(q-1\).
\end{itemize}

Quindi ciascun blocco di Jordan associato a \(\lambda\) fa perdere esattamente \emph{un’unità} di rango.

\begin{proposizione}
Se \(\lambda\) è autovalore di \(A\) con molteplicità geometrica \(m_g(\lambda)=k\), allora
\begin{equation}
\operatorname{rank}(A-\lambda I)=n-k.
\label{eq:L13_rank_deficiency}
\end{equation}
\end{proposizione}

\begin{proof}
La quantità \(k\) coincide col numero di blocchi di Jordan associati a \(\lambda\).  
Ogni blocco \(J_\lambda^{(q_i)}-\lambda I\) ha rango \(q_i-1\), mentre tutti gli altri blocchi hanno rango pieno. Sommando i contributi si ottiene una perdita totale di rango pari a \(k\), cioè
\[
\operatorname{rank}(A-\lambda I)=n-k.
\]
\end{proof}

\begin{osservazione}
La perdita di rango di \(A-\lambda I\) non dipende dalla molteplicità algebrica di \(\lambda\), ma dalla sua \textbf{molteplicità geometrica}, cioè dal numero di blocchi di Jordan associati a \(\lambda\).
\end{osservazione}

\section{Conseguenza importante: numero minimo di ingressi}

Supponiamo che \(\lambda\) sia un autovalore con
\[
m_g(\lambda)=k.
\]
Allora \(A-\lambda I\) ha rango \(n-k\). Per fare in modo che la matrice aumentata
\[
\begin{bmatrix}
A-\lambda I & B
\end{bmatrix}
\]
raggiunga rango \(n\), bisogna recuperare almeno \(k\) direzioni indipendenti.

Questo implica il seguente fatto molto utile.

\begin{proposizione}
Se esiste un autovalore \(\lambda\) di \(A\) con molteplicità geometrica \(k\), allora per poter sperare che il sistema sia completamente raggiungibile è necessario avere almeno \(k\) ingressi indipendenti.
\end{proposizione}

\begin{proof}
Ogni colonna di \(B\) può aumentare il rango della matrice aumentata di al più una unità.  
Se \(A-\lambda I\) parte da rango \(n-k\), per arrivare a rango \(n\) occorre recuperare almeno \(k\) unità di rango, quindi servono almeno \(k\) colonne efficaci di \(B\).
\end{proof}

\begin{osservazione}
In particolare, con \textbf{un solo ingresso} è impossibile avere completa raggiungibilità se esiste un autovalore con molteplicità geometrica maggiore di \(1\).
\end{osservazione}

\section{Esempio 1: sistema già in forma standard di raggiungibilità}

Consideriamo
\begin{equation}
A=
\begin{bmatrix}
-1 & 0 & 2\\
0 & 3 & 1\\
0 & 0 & 2
\end{bmatrix},
\qquad
B=
\begin{bmatrix}
1\\
0\\
0
\end{bmatrix}.
\label{eq:L13_es1}
\end{equation}

La matrice \(A\) è già triangolare superiore, e anzi è già nella forma standard di raggiungibilità: l’ingresso agisce soltanto sul primo stato, mentre gli autovalori \(3\) e \(2\) stanno nella parte non raggiungibile.

\subsection*{Test PBH sugli autovalori}

Gli autovalori sono
\[
\lambda_1=-1,\qquad \lambda_2=3,\qquad \lambda_3=2.
\]

\paragraph{Per \(\lambda=-1\):}
\[
\begin{bmatrix}
A+I & B
\end{bmatrix}
=
\begin{bmatrix}
0 & 0 & 2 & 1\\
0 & 4 & 1 & 0\\
0 & 0 & 3 & 0
\end{bmatrix}.
\]
Questa matrice ha rango \(3\).

\paragraph{Per \(\lambda=3\):}
\[
\begin{bmatrix}
A-3I & B
\end{bmatrix}
=
\begin{bmatrix}
-4 & 0 & 2 & 1\\
0 & 0 & 1 & 0\\
0 & 0 & -1 & 0
\end{bmatrix}.
\]
Questa matrice ha rango \(2\).

\paragraph{Per \(\lambda=2\):}
\[
\begin{bmatrix}
A-2I & B
\end{bmatrix}
=
\begin{bmatrix}
-3 & 0 & 2 & 1\\
0 & 1 & 1 & 0\\
0 & 0 & 0 & 0
\end{bmatrix}.
\]
Questa matrice ha rango \(2\).

\subsection*{Conclusione}

Il criterio PBH fallisce per \(\lambda=3\) e per \(\lambda=2\).  
Dunque il sistema \eqref{eq:L13_es1} \textbf{non è completamente raggiungibile}.

Più precisamente:
\begin{itemize}
    \item il modo associato a \(-1\) è raggiungibile;
    \item i modi associati a \(3\) e \(2\) non sono raggiungibili.
\end{itemize}

\section{Esempio 2: due blocchi di Jordan associati allo stesso autovalore}

Consideriamo ora
\begin{equation}
A=
\begin{bmatrix}
-1 & 1 & 0\\
0 & -1 & 0\\
0 & 0 & -1
\end{bmatrix},
\qquad
B_1=
\begin{bmatrix}
0\\
0\\
1
\end{bmatrix}.
\label{eq:L13_es2}
\end{equation}

La matrice \(A\) ha un autovalore unico \(-1\), con:
\[
m_a(-1)=3,
\qquad
m_g(-1)=2.
\]
Infatti i blocchi di Jordan associati a \(-1\) sono due: uno di ordine \(2\) e uno di ordine \(1\).

\subsection*{Test PBH}

Poiché l’unico autovalore è \(-1\), basta controllare
\[
\begin{bmatrix}
A+I & B_1
\end{bmatrix}
=
\begin{bmatrix}
0 & 1 & 0 & 0\\
0 & 0 & 0 & 0\\
0 & 0 & 0 & 1
\end{bmatrix}.
\]
Il rango è
\[
\operatorname{rank}
\begin{bmatrix}
A+I & B_1
\end{bmatrix}
=2<3.
\]
Quindi il sistema non è completamente raggiungibile.

\subsection*{Interpretazione}

Qui il problema non è l’autovalore in sé, ma il fatto che \(-1\) ha due blocchi di Jordan distinti, cioè due catene indipendenti. Con un solo ingresso non si riesce a pilotare entrambe.

\begin{osservazione}
Questo esempio mostra bene la differenza tra:
\begin{itemize}
    \item \emph{molteplicità algebrica} dell’autovalore;
    \item \emph{numero di catene modali indipendenti} associate a quell’autovalore.
\end{itemize}
Nel criterio PBH conta la seconda, cioè la molteplicità geometrica.
\end{osservazione}

\section{Esempio 3: scegliere un altro ingresso non basta necessariamente}

Riprendiamo la stessa matrice \(A\), ma scegliamo ora
\[
B_2=
\begin{bmatrix}
0\\
1\\
0
\end{bmatrix}.
\]
Allora
\[
\begin{bmatrix}
A+I & B_2
\end{bmatrix}
=
\begin{bmatrix}
0 & 1 & 0 & 0\\
0 & 0 & 0 & 1\\
0 & 0 & 0 & 0
\end{bmatrix}.
\]
Anche in questo caso
\[
\operatorname{rank}
\begin{bmatrix}
A+I & B_2
\end{bmatrix}
=2<3.
\]
Quindi il sistema continua a non essere completamente raggiungibile.

\subsection*{Commento}

Con \(B_2\) si riesce a influenzare una catena diversa, ma rimane comunque una direzione modale fuori portata.  
Il punto cruciale è che con un solo ingresso non si possono recuperare le due unità di rango perse da \(A+I\).

\section{Esempio 4: un caso completamente raggiungibile}

Consideriamo ora
\begin{equation}
A=
\begin{bmatrix}
-1 & 0 & 2\\
0 & 3 & 1\\
0 & 0 & 2
\end{bmatrix},
\qquad
B=
\begin{bmatrix}
1\\
0\\
1
\end{bmatrix}.
\label{eq:L13_es4}
\end{equation}

Gli autovalori sono ancora
\[
-1,\quad 3,\quad 2.
\]
Ma il diverso ingresso cambia radicalmente la situazione.

\paragraph{Per \(\lambda=-1\):}
\[
\begin{bmatrix}
A+I & B
\end{bmatrix}
=
\begin{bmatrix}
0 & 0 & 2 & 1\\
0 & 4 & 1 & 0\\
0 & 0 & 3 & 1
\end{bmatrix},
\qquad
\operatorname{rank}=3.
\]

\paragraph{Per \(\lambda=3\):}
\[
\begin{bmatrix}
A-3I & B
\end{bmatrix}
=
\begin{bmatrix}
-4 & 0 & 2 & 1\\
0 & 0 & 1 & 0\\
0 & 0 & -1 & 1
\end{bmatrix},
\qquad
\operatorname{rank}=3.
\]

\paragraph{Per \(\lambda=2\):}
\[
\begin{bmatrix}
A-2I & B
\end{bmatrix}
=
\begin{bmatrix}
-3 & 0 & 2 & 1\\
0 & 1 & 1 & 0\\
0 & 0 & 0 & 1
\end{bmatrix},
\qquad
\operatorname{rank}=3.
\]

\subsection*{Conclusione}

Il criterio PBH è soddisfatto per tutti gli autovalori, quindi il sistema \eqref{eq:L13_es4} è \textbf{completamente raggiungibile}.

\begin{osservazione}
Questo esempio fa capire che non conta solo la presenza di certi autovalori, ma \emph{come} l’ingresso è collegato al sistema.
\end{osservazione}

\section{Interpretazione tramite autovettori sinistri}

Il criterio PBH può essere letto anche così:

\begin{quote}
\emph{il sistema è completamente raggiungibile se non esiste alcun autovettore sinistro di \(A\) ortogonale a tutte le colonne di \(B\).}
\end{quote}

Infatti, come visto nella dimostrazione:
\[
w^\top(A-\lambda I)=0,
\qquad
w^\top B=0
\]
significa contemporaneamente che:
\begin{itemize}
    \item \(w^\top\) è un autovettore sinistro di \(A\) associato a \(\lambda\);
    \item la direzione modale corrispondente non viene eccitata dall’ingresso.
\end{itemize}

Questa è forse la lettura più concettuale dell’intero lemma PBH.

\section{Messaggio pratico della lezione}

Il criterio PBH è potentissimo perché evita, in molti casi, di calcolare esplicitamente la matrice di raggiungibilità
\[
R=\begin{bmatrix}
B & AB & \cdots & A^{n-1}B
\end{bmatrix}.
\]

In pratica:
\begin{itemize}
    \item si trovano gli autovalori di \(A\);
    \item per ciascun autovalore \(\lambda\), si calcola il rango di
    \[
    \begin{bmatrix}
    A-\lambda I & B
    \end{bmatrix};
    \]
    \item se il rango è sempre \(n\), il sistema è completamente raggiungibile;
    \item se per almeno un autovalore il rango scende, allora esiste almeno un modo non raggiungibile.
\end{itemize}

\section{Schema riassuntivo}

\begin{center}
\fbox{
\parbox{0.93\textwidth}{
\textbf{Criterio PBH per la raggiungibilità}

\vspace{0.3cm}

La coppia \((A,B)\) è completamente raggiungibile se e solo se
\[
\operatorname{rank}
\begin{bmatrix}
A-\lambda I & B
\end{bmatrix}=n
\qquad \forall \lambda\in\C.
\]

\begin{itemize}
    \item Se \(\lambda\notin\sigma(A)\), il rango è automaticamente massimo.
    \item I valori davvero critici sono quindi solo gli autovalori di \(A\).
    \item Se \(m_g(\lambda)=k\), allora \(A-\lambda I\) perde \(k\) unità di rango.
    \item Per recuperarle, \(B\) deve fornire almeno \(k\) direzioni indipendenti.
    \item Se esiste un autovettore sinistro \(w^\top\neq 0\) tale che
    \[
    w^\top A=\lambda w^\top,
    \qquad
    w^\top B=0,
    \]
    allora il modo associato a \(\lambda\) non è raggiungibile.
\end{itemize}
}}
\end{center}

\section{Conclusione}

Questa lezione chiude idealmente il discorso iniziato con la forma standard di raggiungibilità:
\begin{itemize}
    \item nella lezione precedente abbiamo separato \emph{geometricamente} la parte raggiungibile da quella non raggiungibile;
    \item in questa lezione abbiamo ottenuto un test \emph{algebrico} molto rapido per riconoscere la stessa proprietà.
\end{itemize}

Il criterio PBH è uno degli strumenti più importanti dell’intero corso, perché mette in relazione:
\begin{itemize}
    \item autovalori;
    \item blocchi di Jordan;
    \item ingressi;
    \item modi effettivamente pilotabili.
\end{itemize}

Nelle lezioni successive questo criterio tornerà continuamente, sia nella progettazione dei controllori sia nell’analisi strutturale dei sistemi.